\chapter{Contribuições Esperadas e Limites da Proposta}

\section{Contribuições esperadas}

Espera-se produzir uma proposta de ensino introdutório de \textit{machine learning} situada no campo da Computação na Educação Básica. O trabalho articula uma plataforma local a uma sequência de oito encontros com orientações de mediação docente, além de instrumentos individuais analisados por critérios próprios. Esse conjunto permitirá examinar como os estudantes ligam os dados às regras e aos parâmetros na saída do modelo, explicando seus limites.

A proposta oferecerá aos estudantes iniciantes um percurso que dispensa programação prévia ou formalização matemática avançada. A expectativa é criar condições para que a turma levante hipóteses antes de testar alterações e explicar o efeito observado nas tabelas e nos gráficos. Mais do que atestar domínio técnico profundo dos modelos, a pesquisa busca identificar compreensões conceituais situadas ao longo da intervenção.

A matriz curricular construída para o projeto explicitará o vínculo dos objetivos de aprendizagem e das atividades propostas com as habilidades de Computação do Ensino Médio que articulam análise de dados e modelos simples aos fundamentos de inteligência artificial. Esse desenho busca contribuir para o debate sobre a Computação na Educação Básica, sem supor que a BNCC prescreva algoritmos específicos ou uma organização curricular única.

A investigação também poderá sistematizar uma experiência docente adaptável a outras turmas de Ensino Médio, respeitando a infraestrutura disponível e a mediação necessária frente ao repertório e ao currículo de cada escola. Longe de constituírem solução pronta ou universal, a plataforma e os materiais oferecem subsídios a professores interessados em abordar inteligência artificial e dados para além do uso passivo de ferramentas prontas.

O produto também tornará visível o planejamento que sustenta sua utilização, já que a simples oferta de telas com gráficos e controles seria insuficiente. Por isso, tanto os cenários da sequência quanto as questões de previsão e os critérios de análise indicarão a compreensão esperada e as explicações observáveis em cada atividade.

Essa documentação oferecerá aos professores uma referência para justificar adaptações da proposta, evitando padronizar aulas em contextos distintos. Nessas adaptações, caberá preservar o vínculo do conteúdo de Computação com a Matemática mobilizada sob mediação docente, impedindo que o trabalho se reduza à simples demonstração de ferramentas.

\section{Limites e riscos reconhecidos}

O número de encontros e a duração prevista da intervenção limitam a profundidade com que cada modelo poderá ser abordado. Trabalhar cinco modelos introdutórios em oito aulas exige retomadas e comparações constantes, mas não permite supor que uma única exploração resulte em aprendizagem consolidada. O projeto deverá priorizar conceitos e relações fundamentais, deixando de fora formalizações e procedimentos técnicos que ultrapassem o nível introdutório definido.

A participação dos estudantes pode refletir diferenças tanto no repertório matemático quanto na familiaridade e no acesso aos computadores. Embora a sequência preveja estratégias de acesso e mediação, tais cuidados estão longe de eliminar essas desigualdades prévias. Os registros docentes ajudarão a ler como essas condições interferem nas produções individuais, evitando tomar o modelo ou a plataforma como explicação exclusiva para as respostas observadas.

O estudo ocorrerá em uma turma sob docência do próprio pesquisador, ainda a ser definida. Assim, as evidências produzidas refletirão a mediação realizada naquele contexto específico, considerando tanto os dados sintéticos escolhidos quanto as condições materiais do campo. Afasta-se a pretensão de generalizar resultados para toda a rede pública ou de medir um efeito causal isolado da plataforma.

Comparar o núcleo inicial ao final permitirá examinar mudanças e permanências nas explicações, enquanto o desafio de transferência trará evidências sobre o uso dos conceitos em nova situação. Ainda assim, tais achados são insuficientes para atribuir à plataforma a causa isolada das mudanças.

Há limites próprios dos modelos e das representações adotadas. Os conjuntos de dados serão simplificados para evidenciar relações, evitando confundir as saídas da interface com previsões confiáveis sobre pessoas ou situações reais. Além disso, as vistas tridimensionais exibirão apenas métricas calculadas, sem acrescentar novos atributos aos dados.

Os elementos didáticos serão revistos antes da aplicação para afastar a ideia de neutralidade ou infalibilidade em sistemas tidos como autônomos na inteligência artificial.
